\subsection{Small-Lesion Segmentation as a Sparse-Signal Learning Problem}
\label{sec:problem}

\paragraph{Overview in plain terms.}
A segmentation network is trained by repeatedly comparing its output with
the reference mask and adjusting itself to reduce the disagreement, measured
by a \emph{loss}. The standard loss counts disagreement voxel by voxel, so a
lesion influences training in proportion to its volume, and a missed lesion
of a few voxels barely changes the loss. CATMIL adds two terms with
different roles. A detection term asks, for every reference lesion, only
whether the network marked at least one of its voxels, and penalizes it
strongly when it did not. A balancing term gives every lesion a comparable
weight in the overlap loss, whatever its volume, so that detected small
lesions are also delineated. The network itself and its use after training
are unchanged.

\paragraph{Notation.}
Let $\Omega \subset \mathbb{Z}^{3}$ denote the voxel domain of a 3D MRI
volume. For each voxel $i \in \Omega$, let $g_i \in \{0,1\}$ be the
ground-truth label and $p_i \in [0,1]$ the foreground probability predicted
by the segmentation network $f_{\theta}$. The set of lesion voxels is
$G = \{ i \in \Omega \mid g_i = 1 \}$, and it decomposes into $K$ disjoint
connected components, each corresponding to one lesion:
\begin{equation}
G = \bigcup_{k=1}^{K} C_k,
\qquad
C_k \cap C_j = \varnothing \quad \text{for } k \neq j ,
\end{equation}
where $|C_k|$ is the number of voxels in lesion $C_k$. In multiple
sclerosis the lesion signal is sparse in two ways. Lesions occupy a very
small part of the volume, $|G| \ll |\Omega|$, and their sizes are highly
unequal: in a survey based on 5-mm-thick slices, lesions with an in-plane
area below 10~mm$^2$ (equivalent diameter below 3.5~mm) accounted for
nearly 20\% of all lesions but for about 1\% of the total lesion load on
average~\cite{wang_survey_1997}.

A voxel-wise loss such as cross-entropy is a sum of per-voxel terms
$\ell(p_i, g_i)$, which can be grouped by lesion:
\begin{equation}
\mathcal{L}_{\mathrm{vox}}
=
\sum_{i \in \Omega} \ell(p_i, g_i)
=
\sum_{k=1}^{K} \sum_{i \in C_k} \ell(p_i, 1)
+
\sum_{i \notin G} \ell(p_i, 0) .
\label{eq:sl_voxel_loss}
\end{equation}
Lesion $C_k$ enters the loss through $|C_k|$ terms, so its influence on
training is proportional to its size, and a lesion of a few voxels barely
changes the objective~\cite{kofler_blob_2023}. Overlap-based losses behave
in the same way and are additionally unstable on small
objects~\cite{azad_loss_2023}. The DSC shows the same insensitivity: if a
prediction is perfect except that one lesion carrying 1\% of the lesion
volume is missed completely, the DSC is still about 0.995. Voxel-level
overlap therefore does not control lesion-level detection.

The training objective should therefore meet two separate requirements.
The first concerns \emph{detection}: the complete omission of a lesion
should be penalized directly, at the level of the lesion rather than of its
voxels. The second concerns \emph{balance}: the weight of lesion $C_k$ in
the loss should not scale with $|C_k|$. Both are imposed on the loss only, so
that differences in performance can be attributed to the objective.

\subsection{CATMIL Objective}

CATMIL adds two terms to the standard nnU-Net loss, one for each
requirement: a lesion-level MIL term for detection and a Component-Adaptive
Tversky (CAT) term to balance lesions. The ablation attributes the detection gain to the MIL term
(\cref{appendix:catmil-components}).
The connected components
$C_1,\ldots,C_K$ are computed from the ground-truth mask of each training
patch using 6-connectivity, the same lesion definition as in the
evaluation.

\subsubsection{Lesion-Level Multiple Instance Learning Term}

Overlap-based losses, including those that balance lesions of different
sizes, do not distinguish a partially segmented lesion from one that is not
detected at all. The detection requirement is addressed with Multiple
Instance Learning~\cite{ilse_attention-based_2018}: each lesion is treated
as a bag of voxels that counts as detected if at least one of its voxels
receives a high foreground probability. The detection score of lesion $C_k$ is its
strongest response, $s_k = \max_{i \in C_k} p_i$, and the MIL term is
\begin{equation}
\mathcal{L}_{\mathrm{MIL}}
=
-\frac{1}{K} \sum_{k=1}^{K} \log \max(s_k, \varepsilon) ,
\label{eq:sl_mil_loss}
\end{equation}
where $\varepsilon=10^{-6}$ keeps each contribution finite and $K$ is the
number of lesions in the training batch.

If a lesion is missed, all its probabilities are close to zero, so
$s_k \approx 0$ and its contribution is large; once a single voxel inside
the lesion is marked, the contribution vanishes. The training signal does
not depend on lesion size and is largest for lesions that are currently
missed (\cref{appendix:catmil-properties}), so a lesion of a few voxels
receives the same detection signal as a large one. The term does not
require the whole lesion to be segmented; delineation is left to the
overlap terms of the base loss.

\subsubsection{Component-Adaptive Tversky Term}

The MIL term asks only whether a lesion is found; once a lesion is
detected, its delineation is again supervised by the voxel-wise base loss,
in which large lesions dominate. The CAT term is designed to remove this
dependence on lesion size in \cref{eq:sl_voxel_loss} by weighting each
voxel according to the lesion it belongs to:
\begin{equation}
w_i =
\begin{cases}
(|C_k| + \varepsilon_{c})^{-1}, & \text{if } i \in C_k, \\
w_{\mathrm{bg}}, & \text{if } g_i = 0 .
\end{cases}
\label{eq:sl_cat_weights}
\end{equation}
The total weight of lesion $C_k$ is then
\begin{equation}
\rho_k
=
\sum_{i \in C_k} w_i
=
\frac{|C_k|}{|C_k| + \varepsilon_{c}} ,
\label{eq:sl_cat_lesion_weight}
\end{equation}
which is bounded by one and no longer grows with lesion volume: with
$\varepsilon_{c}=5$ voxels, it is $0.67$ for a 10-voxel lesion and $0.95$
for a 100-voxel lesion, so lesions whose volumes differ by orders of
magnitude receive weights that differ by less than a factor of six. The
background weight is $w_{\mathrm{bg}}=0.01$. With these weights, the
weighted true-positive, false-positive, and false-negative sums
$\mathrm{TP}_{w}$, $\mathrm{FP}_{w}$, and $\mathrm{FN}_{w}$ enter the
Tversky index~\cite{salehi_tversky_2017}:
\begin{equation}
\mathcal{L}_{\mathrm{CAT}}
=
1 -
\frac{\mathrm{TP}_{w}+\delta}
{\mathrm{TP}_{w} + \alpha\, \mathrm{FP}_{w} + \beta\, \mathrm{FN}_{w} + \delta},
\label{eq:sl_cat_loss}
\end{equation}
with $\alpha=0.3$ and $\beta=0.7$ as in the Tversky baseline. In effect,
the CAT term is a Tversky loss in which the lesion, rather than the voxel,
is the unit of counting; the complete definitions, the remaining constants,
and the resulting gradient are given in \cref{appendix:catmil-properties}.

\subsubsection{Objective Function}

The complete objective adds the two terms to the default nnU-Net loss,
$\mathcal{L}_{\mathrm{base}} = \mathcal{L}_{\mathrm{Dice}} + \mathcal{L}_{\mathrm{CE}}$:
\begin{equation}
\mathcal{L}_{\mathrm{CATMIL}}
=
\mathcal{L}_{\mathrm{base}}
+
\lambda_{\mathrm{MIL}}\, \mathcal{L}_{\mathrm{MIL}}
+
\lambda_{\mathrm{CAT}}(t)\, \mathcal{L}_{\mathrm{CAT}} .
\label{eq:sl_catmil_loss}
\end{equation}
The base loss provides voxel-level classification and overall overlap, the
MIL term lesion-level detection, and the CAT term balanced overlap per
lesion. The MIL weight is
constant, $\lambda_{\mathrm{MIL}}=0.2$. The CAT weight rises linearly from zero to $\lambda_{\mathrm{CAT}}=0.1$ over
the first 50 epochs, so that the model first learns basic lesion--background
separation before the supervision shifts towards lesion balance. Both
weights were fixed before the experiments and used unchanged for both
datasets; their effect is examined in \cref{appendix:catmil-loss-weights}.
